% !TEX root = ../main.tex
% Appendix A: quick guide to running meshsync. Commands: AGENTS.md §3, PROJECT_CONTEXT.md §14,
% scripts/run_local_mesh.sh, scripts/gen_certs.sh. The status output in "What you should see" is
% copied from a real run of scripts/run_local_mesh.sh -n 3 (2026-11-24); the folder path differs.
% Trimmed to one page on 2026-11-24 (Files table dropped; Linux-only disclaimer added).
\section[\texorpdfstring{\protect\fbox{\strut Quick Guide: Running meshsync}\protect\hspace{6pt}%
 \protect\colorbox{black}{\protect\color{white}\protect\bfseries\ START HERE\ }}{Quick Guide: Running meshsync}]{Quick Guide: Running meshsync}
\label{app:run}

meshsync runs on Linux with Python 3.11 or newer and needs no package beyond the standard library.
\textbf{The whole project was built and tested on Linux machines only;} we have not tried it on
Windows or macOS. Run all commands from the repository root (Appendix~\ref{app:code}), after
installing:
\begin{Verbatim}[frame=single,fontsize=\small]
python -m venv .venv && . .venv/bin/activate && pip install -e .[dev]
\end{Verbatim}

\subsection{One machine, three peers}
\begin{Verbatim}[frame=single,fontsize=\small]
scripts/run_local_mesh.sh -n 3
\end{Verbatim}
Add \code{--tls} for TLS 1.3; Ctrl-C stops all peers. Each peer gets its own folder,
\code{.localmesh/\allowbreak peerN/data}, and TCP port 4700N; all share UDP discovery port 47000. In a second
terminal, with the virtual environment active, put a file in one folder and ask another peer:
\begin{Verbatim}[frame=single,fontsize=\small]
echo hello > .localmesh/peer1/data/hello.txt
meshsync status --dir .localmesh/peer2/data
\end{Verbatim}

\subsection{Three laptops}
The hosts must be on one subnet, and the firewall must allow UDP 47000 and TCP 4700x. Make the
secret and the certificates once and copy them to every laptop (keep the key private); then start
each laptop with its own \code{--name}:
\begin{Verbatim}[frame=single,fontsize=\small]
mkdir -p ~/.meshsync ~/MeshSync
python -c 'import secrets; print(secrets.token_hex(32))' > ~/.meshsync/key
chmod 600 ~/.meshsync/key
scripts/gen_certs.sh ~/.meshsync/certs ahmed mujtaba huzaifa
meshsync run --dir ~/MeshSync --name ahmed --secret-file ~/.meshsync/key \
    --tls --cert-dir ~/.meshsync/certs
\end{Verbatim}
Use \code{--plain} instead of \code{--tls} to run without TLS. If broadcast is blocked, add
\code{--peer HOST:PORT} for each other laptop.

\subsection{What you should see}
\nopagebreak\code{meshsync status} on peer2, a few seconds later (real output):
\begin{Verbatim}[frame=single,fontsize=\small]
running
node     peer2 (7F44DFAD), plain, tcp 47002, udp 47000, pid 22653
started  2026-11-24 19:51:15
updated  2026-11-24 19:51:25 (1s ago)
index    rev 1, 1 files, 0 tombstones
peers    2
  NAME   PEER_ID   ADDRESS              STATE   LAST SEEN
  peer3  7E49D2F1  10.176.187.28:47003  SYNCED  11s ago
  peer1  AB2DB2AD  127.0.0.1:47001      SYNCED  5s ago
\end{Verbatim}
\code{hello.txt} is now in all three \code{data} folders.
